\documentclass[11pt]{article}

\usepackage[T1]{fontenc}
\usepackage{amsmath,amssymb,amsthm}
\usepackage{booktabs}
\usepackage{hyperref}
\usepackage[margin=1in]{geometry}
\hypersetup{hidelinks}

\newtheorem{theorem}{Theorem}
\newtheorem{corollary}[theorem]{Corollary}
\newtheorem{lemma}[theorem]{Lemma}
\newtheorem{remark}[theorem]{Remark}
\newcommand{\A}{\mathbf A}
\newcommand{\C}{\mathbf C}

\title{A Cube-Component Exit for Three-Dimensional Keller Maps}
\author{Author information withheld pending priority review}
\date{\today}

\begin{document}
\maketitle

\begin{abstract}
Let \(k\) be an algebraically closed field of characteristic zero.  We
establish the following structural constraint: a polynomial
\[
f=C+\ell^3+Q_2+L_1\in k[x,y,z],
\]
where \(\ell\ne0\) is linear and \(Q_2,L_1\) are homogeneous of degrees
two and one, is a polynomial coordinate whenever it has no critical
point.  The proof is an elementary rank classification of the
quadratic form transverse to \(\ell\), and gives a coordinate
automorphism and inverse of degree at most three.  Over \(\C\), a
Keller map of degree \(d\) having such a target-linear combination of
its components therefore reduces
fibrewise to plane Keller maps of degree at most \(3d\).  Using the
reported plane counterexample floor \(108\), this proves invertibility
for \(d\le35\); Moh's classical range gives the conservative fallback
\(d\le33\).  We record only a finite, machine-checked application to a
fixed denominator and make no global Jacobian or denominator-row
claim.
\end{abstract}

\noindent\textbf{Draft freeze (UTC):} 2026-07-26T07:53:50Z.
\par\noindent
\textbf{First public release (UTC):}
\texttt{2026-07-26T08:13:13Z}.

\section{The coordinate theorem}

Recall that \(f\in k[x,y,z]\) is a \emph{coordinate} if there are
\(g,h\in k[x,y,z]\) with \(k[f,g,h]=k[x,y,z]\).  We call \(f\)
\emph{nonsingular} if its three partial derivatives have no common
zero in \(k^3\).

\begin{theorem}[Cube-component coordinate theorem]\label{thm:coordinate}
Let \(k\) be algebraically closed of characteristic zero.  Let
\(\ell\) be a nonzero linear form, let \(Q_2,L_1\) be homogeneous of
degrees two and one, and let \(C\in k\).  If
\[
f=C+\ell^3+Q_2+L_1
\]
is nonsingular, then \(f\) is a coordinate of \(k[x,y,z]\).  It can be
included in a polynomial coordinate system whose forward and inverse
maps both have degree at most three.
\end{theorem}

\begin{proof}
A linear source change sends \(\ell\) to \(x\).  A constant translation
removes \(C\).  With \(Y=(y,z)^t\), write
\begin{equation}\label{eq:normal}
 f=x^3+A x^2+x\,c^tY+\frac12Y^tMY+a x+d^tY,
\end{equation}
where \(M=M^t\in M_2(k)\) and \(c,d\in k^2\).  Then
\begin{equation}\label{eq:gradient}
 \nabla_Yf=MY+cx+d,\qquad
 f_x=3x^2+2Ax+c^tY+a.
\end{equation}
We split by the rank of \(M\).

\smallskip
\noindent\emph{Rank two.}
If \(M\) is invertible, \(\nabla_Yf=0\) has the unique solution
\[
Y=-M^{-1}(cx+d).
\]
After this substitution, \(f_x\) is the explicit univariate quadratic
\begin{equation}\label{eq:ranktworesidual}
3x^2+\bigl(2A-c^tM^{-1}c\bigr)x+
       \bigl(a-c^tM^{-1}d\bigr).
\end{equation}
Its leading coefficient is \(3\), so it has a root \(x_0\in k\).
Taking \(Y_0=-M^{-1}(cx_0+d)\) makes every derivative in
\eqref{eq:gradient} vanish.  Thus this rank is impossible.

\smallskip
\noindent\emph{Rank one.}
By a constant linear change in \(Y\), equation~\eqref{eq:normal} becomes
\begin{equation}\label{eq:rankone}
f=g(x)+x(\beta u+\gamma v)+\frac m2u^2+\varepsilon u+\eta v,
\qquad m\ne0.
\end{equation}
Here \(g(x)=x^3+Ax^2+ax\).  If \(\gamma\ne0\), the explicit point
\begin{equation}\label{eq:rankonewitness}
x_0=-\frac{\eta}{\gamma},\qquad
u_0=-\frac{\beta x_0+\varepsilon}{m},\qquad
v_0=-\frac{g'(x_0)+\beta u_0}{\gamma}
\end{equation}
satisfies \(f_v=f_u=f_x=0\), in that order.  Suppose \(\gamma=0\).
If also \(\eta=0\), put
\[
u(x)=-\frac{\beta x+\varepsilon}{m}.
\]
Then \(f_u=f_v=0\), while
\begin{equation}\label{eq:rankoneflat}
f_x\bigl(x,u(x),v\bigr)
=3x^2+\left(2A-\frac{\beta^2}{m}\right)x
  +\left(a-\frac{\beta\varepsilon}{m}\right).
\end{equation}
Choose a root \(x_0\), set \(u_0=u(x_0)\), and take \(v_0=0\).
This is a critical point.  Nonsingularity therefore forces
\(\gamma=0,\eta\ne0\).  In this last case
\[
f=\eta v+R(x,u),\qquad
R=g(x)+\beta xu+\frac m2u^2+\varepsilon u.
\]
The coordinate map \((x,u,v)\mapsto(x,u,f)\) has inverse
\begin{equation}\label{eq:rankoneinverse}
(x,u,w)\longmapsto
\left(x,u,\frac{w-R(x,u)}{\eta}\right).
\end{equation}

\smallskip
\noindent\emph{Rank zero.}
Write the rank-zero expression with scalar coefficients as
\begin{equation}\label{eq:rankzero}
 f=g(x)+x(by+cz)+ey+hz.
\end{equation}
Put \(\Delta=bh-ce\).  If \(\Delta\ne0\), use
\[
u=by+cz,\qquad v=ey+hz.
\]
Then
\[
f=g(x)+xu+v,\qquad v=f-g(x)-xu,
\]
and recover the original transverse variables by
\begin{equation}\label{eq:rankzerolinearinverse}
y=\frac{hu-cv}{\Delta},\qquad
z=\frac{-eu+bv}{\Delta}.
\end{equation}
Thus \((x,u,f)\) is a coordinate system with a degree-three inverse.

Assume \(\Delta=0\).  First suppose \((b,c)=(0,0)\).  If
\((e,h)\ne(0,0)\), take \(v=ey+hz\) and a complementary coordinate.
When \(e\ne0\), choose \(u=z\), so
\[
z=u,\qquad y=\frac{v-hu}{e};
\]
when \(e=0\), necessarily \(h\ne0\), and one may choose \(u=y\), so
\[
y=u,\qquad z=\frac vh.
\]
In both subcharts
\[
f=g(x)+v,\qquad v=f-g(x).
\]
If also \((e,h)=(0,0)\), choose a root \(x_0\) of \(g'\) and set
\(y_0=z_0=0\); this is a critical point.

Finally suppose \((b,c)\ne(0,0)\).  Since \(\Delta=0\), there is
\(\delta\in k\) with \((e,h)=\delta(b,c)\).  This assertion involves
no hidden division: use \(\delta=e/b\) on \(b\ne0\), and
\(\delta=h/c\) on the boundary \(b=0,c\ne0\).  At \(x_0=-\delta\),
both transverse derivatives vanish.  If \(b\ne0\), take
\[
y_0=-\frac{g'(x_0)}b,\qquad z_0=0;
\]
if \(b=0,c\ne0\), take
\[
y_0=0,\qquad z_0=-\frac{g'(x_0)}c.
\]
Then \(f_x=0\) as well, giving a critical point.

The only nonsingular cases are therefore the explicit coordinate cases
above.  Their formulas, together with the initial linear changes, have
degree at most three in both directions.
\end{proof}

\begin{remark}[Pivot completeness]
For
\(
M=\left(\begin{smallmatrix}m_{11}&m_{12}\\m_{12}&m_{22}\end{smallmatrix}\right)
\),
the rank-one chart \(m_{11}\ne0\) uses the null vector
\((-m_{12},m_{11})\); explicitly, the quadratic part is
\[
\frac{m_{11}}2\left(y+\frac{m_{12}}{m_{11}}z\right)^2.
\]
On the boundary \(m_{11}=0=\det M\), one has
\(m_{12}=0\), and rank one is exactly \(m_{22}\ne0\).  The remaining
boundary is \(M=0\).  In rank zero, \(\det[c\ d]\ne0\) is the
independent chart.  Its zero boundary splits into \(c=0,d\ne0\),
\(c=d=0\), and \(c\ne0,d=\delta c\); the last is covered by the two
ordinary pivots on the components of \(c\).  Thus no coefficient is
divided by without its zero locus being included.
\end{remark}

\section{The Keller exit over \(\C\)}

\begin{lemma}[Fibre reduction]\label{lem:fibre}
Let \(F:\C^3\to\C^3\) have constant nonzero Jacobian and degree \(d\).
Suppose one component of \(F\) is a coordinate belonging to a
coordinate automorphism whose inverse has degree at most \(e\).  If
every complex plane Keller map of degree at most \(ed\) is an
automorphism, then \(F\) is an automorphism.
\end{lemma}

\begin{proof}
Permute the target components and precompose by the coordinate inverse
to obtain
\[
G=(G_1(u,v,w),G_2(u,v,w),w).
\]
Its Jacobian determinant is a nonzero constant, and expansion along the
last row gives
\[
\frac{\partial(G_1,G_2)}{\partial(u,v)}\in\C^\times.
\]
Moreover, \(\deg G_i\le ed\).  For every \(w_0\in\C\), the specialized
pair \((G_1(u,v,w_0),G_2(u,v,w_0))\) is therefore a plane Keller map in
the assumed range and hence an automorphism.  Thus \(G\) is injective
fibre by fibre.  Ax--Grothendieck implies that the injective polynomial
self-map \(G\) of \(\A^3_\C\) is surjective.  The Keller condition makes
\(G\) étale.  An injective étale morphism over an algebraically closed
field is universally injective and hence an open immersion
\cite{StacksEtale}; surjectivity makes that open immersion an
isomorphism.  Thus \(G\), and therefore \(F\), is an automorphism.
\end{proof}

\begin{corollary}[Degree \(35\), using the plane floor \(108\)]
\label{cor:35}
Let \(F:\C^3\to\C^3\) be a Keller map of degree \(d\).  Suppose there is
a nonzero target covector \(\alpha\in(\C^3)^*\) such that
\[
\alpha\!\cdot\!F=C+\ell^3+Q_2+L_1
\]
with \(\ell\ne0\) linear.  If the plane Keller conclusion is valid for
maximum degree \(<108\), then \(F\) is an automorphism for \(d\le35\).
\end{corollary}

\begin{proof}
Extend \(\alpha\) to a matrix \(T\in\mathrm{GL}_3(\C)\), and replace
\(F\) by \(TF\), placing \(\alpha\!\cdot\!F\) in the third component.
The new determinant is
\[
\det J(TF)=(\det T)(\det JF)\in\C^\times.
\]
Consequently the gradient of \(\alpha\!\cdot\!F\), now the third row of
an everywhere invertible Jacobian matrix, is nowhere zero.  Apply
Theorem~\ref{thm:coordinate} with \(e=3\), then
Lemma~\ref{lem:fibre}.  For \(d\le35\),
\[
ed\le3\cdot35=105<108.
\]
\end{proof}

The plane input in Corollary~\ref{cor:35} is the lower floor reported
in~\cite{GGHV}.  We verified the numerical claim in the primary
preprint; a peer-reviewed journal version was not located.

\begin{corollary}[Conservative Moh fallback]
Under the hypotheses of Corollary~\ref{cor:35}, \(F\) is an
automorphism for \(d\le33\) using only the classical plane range
\(<100\).
\end{corollary}

\begin{proof}
Here \(3d\le99<100\), so the same fibre argument applies with
Moh's low-degree result~\cite{Moh}.
\end{proof}

\section{Finite denominator application}

The theorem is applied only to the entries in Table~\ref{tab:bridge}.
The table is machine-checked against a frozen denominator containing
exactly \(26\) families.

\begin{table}[h]
\centering
\begin{tabular}{lll}
\toprule
Identifier & Cube component & Status here\\
\midrule
\texttt{PF-BRANCH-FOURTH-THIRD} & \(R=p^3\) & newly excluded\\
\texttt{D3-BB-30} & \(R=p^3\) & newly excluded\\
\texttt{D3-OB-300} & \(R=p^3\) & newly excluded\\
\texttt{D4-DN-3} & \(R=L^3\) & already excluded\\
\texttt{D3-SF-20C}, \(z=3\) only & \(R=X^3\) &
newly excluded pivot only\\
\bottomrule
\end{tabular}
\caption{Exact scope of the frozen denominator bridge.}
\label{tab:bridge}
\end{table}

The generic \texttt{D3-SF-20C} normal form is
\[
R=X^2((5-3z)p+4rq).
\]
Accordingly, neither that whole family nor its reciprocal \(z=1/3\)
sheet is claimed.  No quartic row, other denominator family, or global
degree bound follows from this finite bridge.  The three newly excluded
whole families move the frozen fine count from \(6/26\) to \(9/26\).
The containing row remains open, the global quartic status remains
\(4/14\) certified, \(4/14\) provisional, and \(6/14\) open, and the
universal total-degree floor remains four.

\section{Reproducibility and priority}

The dependency-free exact verifier checks the critical-point witnesses,
all coordinate inverses, the degree boundaries
\[
3\cdot35=105<108,\quad 3\cdot36=108,\quad
3\cdot33=99<100,
\]
and the exact denominator bridge.  The primary strict wrapper
\texttt{verify\_strict.sh} requires deliberate mutations of each of
these claims to fail and prints
\texttt{CUBE\_COMPONENT\_EXIT\_STRICT\_PASS}.  The aggregate wrapper
\texttt{verify\_all\_strict.sh} additionally requires the independently
implemented hostile audit and prints
\texttt{CUBE\_COMPONENT\_ALL\_STRICT\_PASS} only after both suites pass.

The precise coordinate lemma was not located in the limited literature
search accompanying this draft.  Related work includes classifications
of cubic polynomials at infinity~\cite{Ribeiro}, broad criteria for
polynomial variables~\cite{Kaliman}, Vistoli's theorem for entire
degree-three Keller maps~\cite{Vistoli}, and the classification of
degree-three automorphisms by Blanc and van Santen~\cite{BlancVanSanten}.
None was found to state the present component criterion.  Worldwide
novelty and priority remain unresolved.

\section*{Status and disclosure}

This is an AI-assisted research draft and has not been peer reviewed.
The exact primary and hostile verification suites check the rank atlas,
inverse formulas, degree arithmetic, and frozen-data bridge encoded in
the accompanying scripts.  Those checks are evidence about the encoded
algebra, not peer review, and do not mechanically verify the external
plane-degree theorem, Ax--Grothendieck, or worldwide priority.  The
separate hostile algebra audit and hypothesis/reference audit are
included in the release directory.

\begin{thebibliography}{9}

\bibitem{Ax}
J.~Ax,
\newblock Injective endomorphisms of varieties and schemes,
\newblock \emph{Pacific J. Math.} 31 (1969), 1--7.

\bibitem{BlancVanSanten}
J.~Blanc and I.~van Santen,
\newblock Automorphisms of the affine 3-space of degree 3,
\newblock \emph{Indiana Univ. Math. J.} 71 (2022), 857--912.

\bibitem{GGHV}
J.~A. Guccione, J.~J. Guccione, R.~Horruitiner, and C.~Valqui,
\newblock Increasing the degree of a possible counterexample to the
Jacobian Conjecture from 100 to 108,
\newblock arXiv:2204.14178 (2022),
\newblock \url{https://arxiv.org/abs/2204.14178}.

\bibitem{Kaliman}
S.~Kaliman,
\newblock Polynomials with general \(\C^2\)-fibers are variables. I,
\newblock arXiv:math/9912058,
\newblock \url{https://arxiv.org/abs/math/9912058}.

\bibitem{Moh}
T.~T. Moh,
\newblock On the Jacobian conjecture and the configurations of roots,
\newblock \emph{J. Reine Angew. Math.} 340 (1983), 140--212.

\bibitem{Ribeiro}
N.~R. Ribeiro,
\newblock Classification at infinity of degree-three polynomials in
three variables,
\newblock arXiv:2201.11026,
\newblock \url{https://arxiv.org/abs/2201.11026}.

\bibitem{StacksEtale}
The Stacks Project Authors,
\newblock Étale and universally injective morphisms are open immersions,
\newblock Tag 025G,
\newblock \url{https://stacks.math.columbia.edu/tag/025G}.

\bibitem{Vistoli}
A.~Vistoli,
\newblock The Jacobian conjecture in dimension 3 and degree 3,
\newblock \emph{J. Pure Appl. Algebra} 142 (1999), 79--89.

\end{thebibliography}

\end{document}
